\documentclass[a4paper,12pt]{article}
\usepackage[utf8]{inputenc}
\usepackage[T1]{fontenc}
\usepackage[french]{babel}
\usepackage{amsmath,amsfonts,amsthm,amssymb,amsbsy,amscd,mathrsfs}
\usepackage{lmodern}
\usepackage{fourier}
%\usepackage{times,mathptmx}
%\DeclareSymbolFont{lmodern}{U}{euex}{m}{n}
%\DeclareMathSymbol\infty \mathord{lmodern}{"31}
\usepackage[scaled=0.875]{helvet} % ss
\usepackage{textcomp}
\usepackage{dsfont} % pour les symboles des ensembles
\usepackage[left=2cm,right=2cm,top=1.5cm,bottom=2.5cm]{geometry}
\usepackage{fancyhdr} 
\fancyhf{}
\setlength{\headheight}{26pt}
\setlength{\footskip}{18pt}
\renewcommand {\headrulewidth}{0pt}
\renewcommand {\footrulewidth}{0pt}
%\usepackage{lastpage}
%\usepackage{makeidx}
\usepackage{graphicx} % inclure le fichier graphique de pondichery 2010
\usepackage{hyperxmp}
\usepackage[colorlinks=true,pdfstartview=FitV,linkcolor=blue,citecolor=blue,urlcolor=blue]{hyperref}
\pagestyle{fancy}
\usepackage{array}
\usepackage{tabularx}
\usepackage{multirow}
\usepackage{hhline}
\usepackage[table]{xcolor}
\usepackage{arydshln}%%% lignes en tirets dans un tableau 
\usepackage{mathtools} %pour l'alignement des termes d'une matrice
\usepackage{fltpoint}
\usepackage{xcolor,colortbl}
\usepackage{pstricks,pst-all,pst-plot,pstricks-add,pst-tree,pst-3dplot,pst-func}
\usepackage{ifthen}
\usepackage{calc}
\usepackage{esvect}
\usepackage[np]{numprint}
\setlength{\parindent}{0mm}
\usepackage [alwaysadjust]{paralist}
\setdefaultenum {1.}{a)}{}{}
\hyphenpenalty=10000 
\tolerance=2000
\emergencystretch=1em
\widowpenalty=10000
\clubpenalty=10000
\raggedbottom
\renewcommand*{\tabularxcolumn}[1]{m{#1}} %centrage vertical des cellules d'un tableau tabularx
\newcommand{\Oij}{$\left( {{\mathrm{O}};\vec i,\vec j} \right)$}
\newcommand{\Oijk}{$\left( {{\mathrm{O}};\vec i,\vec j ,\vec k} \right)$}
\newcommand{\euro}{\texteuro{}}
\newcommand{\R}{\mathds {R}}
%\newcommand{\C}{\mathds {C}}
\newcommand{\N}{\mathds {N}}
\newcommand{\Z}{\mathds {Z}}
\newcommand{\Q}{\mathds {Q}}
\newcommand{\e}{\mathrm {e}}
\newcommand{\dd}{\,\mathrm{d}}
\DecimalMathComma
%%% Mise en forme des algorithmes
\usepackage[french,boxed,titlenumbered,lined,longend]{algorithm2e}
  \SetKwIF {Si}{SinonSi}{Sinon}{si}{alors}{sinon\_si}{alors}{fin~si}
 \SetKwFor{Tq}{tant\_que~}{~faire~}{fin~tant\_que}
 \SetKwFor{PourCh}{pour\_chaque }{ faire }{fin pour\_chaque}
 \SetKwInput{Sortie}{Sortie}
\newcommand{\Algocmd}[1]{\textsf{\textsc{\textbf{#1}}}}\SetKwSty{Algocmd}
  \newcommand{\AlgCommentaire}[1]{\textsl{\small  #1}}  
  
  %%% CADRES
\newrgbcolor{Rfond}{.995 0.945 .985}  
\newrgbcolor{Rougef}{.8 0.1 .2}
\newrgbcolor{Rbord}{.6 0.1 .4}
\newrgbcolor{bleu}{0.1 0.05 .5}

\newcommand{\cadreR}[1]{%
\psframebox[framesep=.8em,linecolor=black,linewidth=0.25pt,framearc=0.1,fillstyle=solid,fillcolor=Rfond]{\begin{minipage}{\linewidth-1.6em}%
#1\end{minipage}}}
\newcommand{\cadre}[1]{%
\psframebox[framesep=.8em,linecolor=black,linewidth=0.25pt,framearc=0]{\begin{minipage}{\linewidth-1.6em}%
#1\end{minipage}}}

\newcommand{\encadrer}[1]{{\setlength{\fboxsep}{2mm}\fbox{#1}}}
%%%%%%%%%%%%%%%
\newcounter{nexo}    % déclaration du numéro d'exo
\setcounter{nexo}{0} % initialisation du numero
\newcommand{\exo}{%
  \stepcounter{nexo} 
  \par{\textsf{\textbf{\textcolor{blue}{{\textsc{exercice  \arabic{nexo}}}}}}}
  }%
\newcommand{\QCM}[3]{ $\bullet$ \quad #1 \hfill 
$\bullet$ \quad #2 \hfill 
$\bullet$ \quad #3 }
\newcommand{\QCMc}[4]{   #1 \\ 
\makebox[3cm][l]{\psframe(.25,.25) \qquad #2 }\hfill 
\makebox[3cm][l]{\psframe(.25,.25) \qquad #3 }\hfill 
\makebox[3cm][l]{\psframe(.25,.25) \qquad #4 }} % QCM avec cases à cocher%
% en-tête
\lhead{NOM - PRENOM:} %gauche
\chead{\textsf{\textcolor{red}{\textbf{\textsc{DS03}}}}} %centre
%\rhead{\footnotesize{\fbox{\textsf{S1E2}}}} %droit
%pied de page
%\lfoot{\scriptsize{\textsf{\textsc{A. Yallouz}}}} %gauche
\cfoot{\scriptsize{\textsf{-~\thepage{}~-}}}  %centre                                                   
%\rfoot{\htmladdnormallink{\scriptsize\textsf{\textsc{MATH@ES}}}{http://yallouz.arie.free.fr}} %droit 
\hypersetup{
pdftitle={titre du document.}, 
pdfauthor={Arié Yallouz},   
pdfsubject={Exercices Terminale ES.}, 
pdfkeywords={exercices, sujets.},
pdfcontacturl={http//yallouz.arie.free.fr/},
pdflang={fr},
baseurl={http://yallouz.arie.free.fr/serveurtex/exostex.php}
}
\begin{document}

\exo

 \medskip 

\medskip 


Étude de la production de plats préparés sous vide.

\medskip

\textbf{Les questions 1, 2 et 3 de cet exercice peuvent être traitées de manière
indépendante. Les résultats seront arrondis, si nécessaire, à } \boldmath $10^{-3}$\unboldmath{}.

L'entreprise BUENPLATO produit en grande quantité des plats préparés sous vide.

L'objectif de cet exercice est d'analyser la qualité de cette production en exploitant
divers outils mathématiques.

\medskip

\begin{enumerate}
\item Sur les emballages, il est précisé que la masse des plats préparés est de 

$400$~grammes. Un plat est conforme lorsque sa masse, exprimée en gramme, est
supérieure à $394$~grammes.

On note $M$ la variable aléatoire qui, à chaque plat prélevé au hasard dans la
production, associe sa masse en gramme. On suppose que la variable aléatoire
$M$ suit la loi normale d'espérance $400$ et d'écart type $5$.
	\begin{enumerate}
		\item Déterminer la probabilité qu'un plat prélevé au hasard ait une masse
comprise entre $394$ et $404$ grammes.


\cadre{
\vspace{1cm}\hspace{18cm}
}


		\item  Déterminer la probabilité qu'un plat soit conforme.
		
		\cadre{
\vspace{1cm}\hspace{18cm}
}

 	\end{enumerate}
\item Les plats préparés sont livrés à un supermarché par lot de $300$.
	
On arrondit la probabilité de l'évènement \og un plat préparé prélevé au hasard
dans la production n'est pas conforme \fg{} à $0,12$.
	
On prélève au hasard $300$~plats dans la production. La production est assez
importante pour que lion puisse assimiler ce prélèvement à un tirage aléatoire
avec remise.
	
On considère la variable aléatoire $X$ qui, à un lot de $300$~plats, associe le nombre
de plats préparés non conformes qu'il contient.
	\begin{enumerate}
		\item Justifier que la variable aléatoire $X$ suit une loi binomiale dont on précisera
les paramètres.

\cadre{
\vspace{5cm}\hspace{18cm}
}

		\item Calculer l'espérance mathématique E$(X)$ et en donner une interprétation.
		
		\cadre{
\vspace{2cm}\hspace{18cm}
}

		\item Calculer la probabilité que dans un échantillon de $300$~plats prélevés au
hasard, au moins $35$ plats soient non conformes.

\cadre{
\vspace{3cm}\hspace{18cm}
}
 	\end{enumerate}
 	
 	\end{enumerate}

 
 \bigskip
 
 \exo
 
 \medskip
 
Un revendeur de matériel photographique désire s'implanter dans une galerie marchande.

Il estime qu'il pourra vendre 40 appareils par jour et les ventes sont deux à deux indépendantes.

Une étude lui a montré que, parmi les différentes marques disponibles, la marque A réalise 38.6\% du marché

\begin{enumerate}
	\item On note X la variable aléatoire qui, un jour donné, associe le nombre d'appareils de marque A vendus ce jour-là.
	
	\begin{enumerate}
		\item On suppose que X suit une loi binomiale. Préciser les paramètres de cette loi.
		
		\cadre{
\vspace{1cm}\hspace{18cm}
}
		\item Calculer la probabilité que, sur 40 appareils vendus par jour, 20 soient de la marque A.
		
		\cadre{
\vspace{1cm}\hspace{18cm}
}
		\item Calculer l'espérance de X ainsi que l'écart type de X. 
		
		\cadre{
\vspace{1cm}\hspace{18cm}
}
	\end{enumerate}
	
	\item On décide d'approcher cette loi par une variable aléatoire Y qui suit une loi normale de paramètre $\mu$ et $\sigma$
	
	\begin{enumerate}
		\item Expliquer pourquoi il est possible d'approcher la loi binomiale par la loi normale.
		
		\cadre{
\vspace{1.8cm}\hspace{18cm}
}
		\item Expliquer pourquoi on peut choisir $\mu=15.44$ et $\sigma=3$.
		
		\cadre{
\vspace{1.8cm}\hspace{18cm}
}
		\item Calculer à l'aide de la loi normale la probabilité $P(19\leqslant X\leqslant 21)$
		
		\cadre{
\vspace{1.8cm}\hspace{18cm}
}
	\end{enumerate}
\end{enumerate}


  \end{document}